\section{What a straight-track reconstruction misses}\label{sec:straight}
\sectionstate{partial}{new section (Dylan 2026-11-10: ``a section about the downsides of assuming there's a track, what do we
miss (bending things)''). Data: the $\delta$-ray, deposit and position-tail classes of \texttt{rerun/effbudget/single} and
\texttt{rerun/effbudget/pass1} (the old sections 5.4 and 5.6 and appendix A), the saturated-blob test
\texttt{rerun/satstudy/satblob.py}, the tails of \texttt{figs/out/angle\_resolution\_v2.json}. Toy:
\texttt{papers/paper1\_v2/figs/make\_gas\_scatter.py}, \texttt{out/gas\_scatter.json}. Displays: \texttt{figs/make\_displays\_v2.py}.}
\intent{The reconstruction assumes one straight track per particle. Say what that costs on the data we have, and what it
would cost on the particles the chambers were built for, which the data cannot show.}

The reconstruction of \cref{sec:reco} fits a straight line with a free charge profile along it. Whatever in the
chamber is not a straight line is either fitted as one anyway or left unexplained. On cosmic muons this assumption
is nearly exact for the muon itself and fails for what the muon produces; for the electrons the chambers were built
for it is not exact even for the particle.

\subsection{On muons: what the muon makes}\label{sec:delta}
A GeV muon is straight over 30\,mm of gas to a few thousandths of a degree. What bends is its company
(\cref{fig:notstraight}). The largest loss of single tracks, 2.4\,\% (\cref{sec:single}), is a muon with a knock-on
electron, a $\delta$-ray, and it takes three forms with the size of the deposit. A $\delta$-ray that leaves along the
track adds charge on one side of it and pulls the fitted line by more than 5\,mm (0.81\,\%; a clean muon covers a
median of 12 strips, these 34). A compact $\delta$-ray saturates a few strips and tilts the line towards itself
(0.87\,\%; 19 strips). A large branching one exceeds the 50 strips at which the reconstruction takes the event for a
discharge, and nothing is fitted (0.76\,\%; 67 strips). None of the three comes with extra activity in the reference
planes above or below the chamber, so the deposit is made by the muon in the chamber's own gas, cathode or window; the
rate is of the order expected for knock-on electrons above 0.1\,MeV. A smaller class is a second deposit that is not
shaped like a particle crossing the gap at all (0.18\,\%).

These events are also the tail of the angular resolution. The 3.4\,\% of tracks reconstructed more than 5\textdegree{}
from the reference (\cref{sec:angle}) are fits pulled far steeper than the muon, not a wider core.

The model cannot be patched cheaply. Adding a compact deposit to the forward model, and excluding the saturated
samples from the fit, raises the agreement with the reference on 350 views with a saturated blob from 71.4 to at
most 72.9\,\%: the deposit is itself a short, curling track over 5 to 10 strips, and not a point. Recovering
these events needs a reconstruction that first recognises the topology, and then fits the muon on the charge that
belongs to it. That could return up to the 2.4 points lost here; how much of it is reachable has not been measured.

\begin{figure}[tbp]
  \centering
  \figorbox[width=0.86\linewidth]{not_straight}{four chamber A events that are not one straight track}
  \caption{Four events of chamber~A (Run~1) that are not a single straight track, both views each (strip position
  against drift time; amplitude normalised to the full scale of the ADC; reference track solid black, fit dashed).
  (a) A $\delta$-ray leaving along the track pulls the straight-line fit. (b) A saturated $\delta$-ray tilts it. (c) A
  shower at the crossing: more than 50 strips, taken for a discharge, nothing fitted. (d) A second deposit that is
  not particle-like. The same events appear with the other loss classes in \cref{app:effdef}.}
  \label{fig:notstraight}
\end{figure}

\subsection{On electrons: the particle itself}\label{sec:gasscatter}
The data of this paper are cosmic muons and high-energy beam particles. The chambers were built for electrons and
positrons of 3 to 16\,MeV, and for these the track is not straight even inside the chamber: they scatter in the gas
while they are being measured. No data set of this paper contains such electrons with a known direction, so this
cannot be measured here; a simple simulation gives its size.

\Cref{fig:gasscatter} follows electrons through the 28.2\,mm of drift gas, $2.4\times10^{-4}$ radiation lengths, with
multiple scattering as a Gaussian process of the Highland width. An 8\,MeV electron changes direction by
0.96\textdegree{} rms between entering and leaving the gap, and a 3\,MeV electron by 2.4\textdegree. The straight line
that best fits such a path, sampled uniformly along the gap as the TPC samples it, differs from the direction at
entry by 0.61 of that: 0.59\textdegree{} at 8\,MeV and 1.5\textdegree{} at 3\,MeV.

Two things follow. First, the bend is invisible to the chamber: the path deviates from its fitted line by 0.02 to
0.06\,mm, well below the position resolution, so the fit is as good as for a muon and nothing flags the event. Second,
its size is that of the chamber's own resolution. For leptons below about 5\,MeV the direction of the line in the gas
differs from the direction with which the particle entered by more than the 1\textdegree{} measured on muons
(\crefp{fig:gasscatter}{b}). The resolution of \cref{sec:angle} is thus the detector's term; for the leptons of
interest the gas adds its own, and the foils and the air in front add more (\cref{tab:material}).
\askDylan{Caveats to keep or state: the Highland formula is used below its stated range ($10^{-3}X_0$; a Lynch--Dahl estimate
is 5\,\% lower), single-scatter tails are not modelled, normal incidence only, and the two terms (chamber on muons, scattering
in the gas) are not combined in one curve. A Geant4 check of this exists in principle in the full simulation; it was not run here.}

\begin{figure}[tbp]
  \centering
  \figorbox[width=0.92\linewidth]{gas_scatter}{simulated electron paths in the drift gas and the angle they cost}
  \caption{Multiple scattering in the 28.2\,mm of drift gas (Ar/iC$_4$H$_{10}$ 95/5, $2.4\times10^{-4}$ radiation
  lengths; Gaussian model of the Highland width, no single-scatter tails). (a) Five simulated paths each for a 4\,GeV
  muon and for electrons of 16, 8 and 3\,MeV, the transverse axis stretched 2.9 times, with the straight line fitted to
  one of them. (b) The rms angle between entry and exit (black; dotted, a Lynch--Dahl estimate) and between the fitted
  line and the direction at entry (orange), against the kinetic energy; dashed, the resolution of the chamber measured
  on muons. (c) The rms deviation of the path from its fitted line, with the strip pitch and the position resolution
  of the chamber for scale.}
  \label{fig:gasscatter}
\end{figure}

\subsection{What the model leaves out by construction}
Three further things are outside a one-straight-track model and are not measured in this paper. A low-energy
electron that stops or scatters through large angles inside the gap leaves a deposit that does not span the gap as a
column, and is not sought. A track with a single hard scatter, a kink, is fitted with one line. And any event with more than 50
strips above threshold is set aside as a discharge before a fit is tried: of the tracks so lost, the true discharges
are a minority (0.47\,\% of single tracks, against 0.76\,\% for large $\delta$-rays and 1.0\,\% in the edge zone).
A veto that counts strips per cluster instead of per plane is the first step to getting those back
(\cref{app:moved}).
